% =============================================================
%  Section 5 – Main Results
%  Label: sec:results
% =============================================================

\section{Main Results}\label{sec:results}

The primary finding is that province-level meat-price shocks
generate large, negative household forecast errors.
I present this result first (Section~\ref{sec:eq3}), then
decompose it into the expectation channel
(Section~\ref{sec:eq1}) and the pass-through channel
(Section~\ref{sec:eq2}).
Section~\ref{sec:event_study} visualises the treatment variation
and the pass-through relationship.

% -------------------------------------------------------------
\subsection{Meat Shocks Generate Systematic Forecast Errors}%
\label{sec:eq3}
% -------------------------------------------------------------

\begin{table}[!htbp]
\centering
\caption{Meat Price Shock and Household Forecast Errors}
\label{tab:eq3_forecast_error}
\begin{threeparttable}
\begin{tabular}{lccccc}
\toprule
& (1) & (2) & (3) & (4) & (5) \\
\multicolumn{6}{l}{\textit{Dep.\ var.: Forecast error ($FE_{it}$)}} \\
\midrule
Meat shock & $-2.007^{***}$ & $-2.786$ & $-5.622^{***}$ & $-6.067^{***}$ & $-4.259$ \\
& $(0.252)$ & $(5.903)$ & $(1.182)$ & $(1.145)$ & $(5.662)$ \\[6pt]
Demographics &  & Yes &  &  & Yes \\
Province FE &  &  & Yes &  &  \\
Wave FE &  &  & Yes &  &  \\
Region $\times$ wave FE &  &  &  & Yes & Yes \\
\midrule
Observations & 69,533 & 22,121 & 69,533 & 69,533 & 22,121 \\
$R^2$ & 0.066 & 0.035 & 0.228 & 0.224 & 0.045 \\
\bottomrule
\end{tabular}
\begin{tablenotes}[flushleft]
\footnotesize
\item \textit{Note.} Standard errors clustered at the province level in parentheses.
\item Forecast error is realized province-level CPI inflation minus the household's directional expectation scaled to CPI units.
\item $^{*}p<0.10$, $^{**}p<0.05$, $^{***}p<0.01$.
\end{tablenotes}
\end{threeparttable}
\end{table}


Table~\ref{tab:eq3_forecast_error} reports estimates of
equation~\eqref{eq:eq3}.
The dependent variable is realised province-level headline CPI
inflation minus the household's inflation expectation formed at
the survey wave.
A negative value indicates overprediction: the household expected
more inflation than materialised.

The preferred estimate under region-by-wave fixed effects
is negative, large, and precisely estimated (column~4).
The coefficient is stable across the specification sequence:
it is negative and statistically distinguishable from zero
in every column, from the bivariate specification through to
the saturated model with fixed effects and demographic controls.
Wild cluster bootstrap confirms significance at the $p < 0.001$
level under both province-plus-wave and region-by-wave fixed-effect
structures (Section~\ref{sec:wcb}).

The economic interpretation is direct.
Provinces that experienced larger meat-price increases between
survey waves had households that subsequently overpredicted
inflation more severely.
The overprediction is not driven by a decline in realised
inflation---the pass-through equation below shows that realised
CPI was if anything lower in high-shock provinces---but by an
upward shift in expectations that exceeded whatever headline
inflation actually materialised.

% -------------------------------------------------------------
\subsection{Decomposition: The Expectation Channel}%
\label{sec:eq1}
% -------------------------------------------------------------

\begin{table}[!htbp]
\centering
\caption{Meat Price Shock and Household Inflation Expectations}
\label{tab:eq1_expectations}
\begin{threeparttable}
\footnotesize
\setlength{\tabcolsep}{3pt}
\begin{tabular}{lcccccc}
\toprule
& (1) & (2) & (3) & (4) & (5) & (6) \\
& Bivariate & Demographics & \shortstack{Province\\+ wave} & \shortstack{Region\\$\times$ wave} & \shortstack{Region $\times$ wave\\+ demographics} & \shortstack{Province +\\region $\times$ wave} \\
\midrule
Meat shock & $0.028^{**}$ & $2.492^{***}$ & $0.090^{*}$ & $0.242$ & $2.636^{***}$ & $0.119^{**}$ \\
& $(0.012)$ & $(0.812)$ & $(0.053)$ & $(0.174)$ & $(0.806)$ & $(0.050)$ \\[6pt]
Demographics &  & Yes &  &  & Yes &  \\
Province FE &  &  & Yes &  &  & Yes \\
Wave FE &  &  & Yes &  &  &  \\
Region $\times$ wave FE &  &  &  & Yes & Yes & Yes \\
\midrule
Observations & 90,133 & 22,121 & 90,133 & 90,133 & 22,121 & 90,133 \\
$R^2$ & 0.000 & 0.038 & 0.011 & 0.002 & 0.042 & 0.012 \\
\bottomrule
\end{tabular}
\begin{tablenotes}[flushleft]
\footnotesize
\item \textit{Note.} The dependent variable is the household's ordinal inflation expectation ($\mu_{it} \in \{-1, 0, 1\}$). Standard errors clustered at the province level in parentheses. Column~(6) is the preferred specification.
\item $^{*}p<0.10$, $^{**}p<0.05$, $^{***}p<0.01$.
\end{tablenotes}
\end{threeparttable}
\end{table}


Table~\ref{tab:eq1_expectations} decomposes the forecast-error
result by asking whether meat-price shocks shift the expectation
itself (equation~\eqref{eq:eq1}).

The coefficient is positive across all five columns.
In the bivariate specification, the estimate is small and
statistically significant at conventional levels.
Under region-by-wave fixed effects (column~4), the point estimate
increases but precision falls as the fixed effects absorb a
substantial share of the province-level shock variation.
With both fixed effects and demographics (column~5), the point
estimate is larger and statistically significant, though the sample
shrinks by approximately three-quarters due to missing demographic
variables.

The positive sign is economically meaningful even where statistical
precision is limited.
Meat's official CPI weight is approximately 5 percent, so the
rational expectation response to a unit meat shock is
$\beta_1^{\text{rational}} \approx 0.05$.
The estimated response in the preferred specification is several
times this benchmark.
A confidence interval that includes zero on this coefficient
does not imply that expectations do not respond; it implies that
the response is imprecisely measured in the demanding FE structure,
which absorbs region-level common shocks.
The forecast-error result in Table~\ref{tab:eq3_forecast_error}
provides the more powerful test: even if the individual expectation
channel is noisy, the joint outcome---expectations minus
realised CPI---is precisely estimated as negative.

% -------------------------------------------------------------
\subsection{Decomposition: The Pass-Through Channel}%
\label{sec:eq2}
% -------------------------------------------------------------

\begin{table}[!htbp]
\centering
\caption{Meat Price Shock and Realized Headline CPI}
\label{tab:eq2_cpi_passthrough}
\begin{threeparttable}
\begin{tabular}{lccc}
\toprule
& (1) & (2) & (3) \\
\multicolumn{4}{l}{\textit{Dep.\ var.: Realized headline CPI (\%, annualized)}} \\
\midrule
Meat shock & $-1.812^{***}$ & $-4.703^{***}$ & $-3.519^{*}$ \\
& $(0.232)$ & $(1.646)$ & $(1.954)$ \\[6pt]
Province FE &  & Yes &  \\
Wave FE &  & Yes &  \\
Region $\times$ wave FE &  &  & Yes \\
\midrule
Observations & 93 & 93 & 93 \\
$R^2$ & 0.192 & 0.704 & 0.589 \\
\bottomrule
\end{tabular}
\begin{tablenotes}[flushleft]
\footnotesize
\item \textit{Note.} Unit of observation is province $\times$ wave.
\item Standard errors clustered at the province level.
\item Dependent variable is annualized realized headline CPI inflation in the province over the subsequent inter-wave period.
\item $^{*}p<0.10$, $^{**}p<0.05$, $^{***}p<0.01$.
\end{tablenotes}
\end{threeparttable}
\end{table}


Table~\ref{tab:eq2_cpi_passthrough} estimates equation~\eqref{eq:eq2}
at the province level.
The dependent variable is forward realised headline CPI inflation
over the subsequent inter-wave period.

The point estimate is negative across all columns: provinces with
larger meat-price shocks experienced slightly lower forward
headline inflation.
This is consistent with the supply-shock reversal documented in
other commodity episodes---the same ASF shock that raised meat
prices in 2019--2020 suppressed them in 2021--2022 as herds
recovered and overproduction set in.

The economic content of this finding is that meat-price shocks do
not predict higher future headline inflation at the province level.
The absence of positive pass-through closes the second half of
the overreaction logic.
Households raise their expectations (Table~\ref{tab:eq1_expectations})
in response to a shock that does not, on net, raise the aggregate
price level they are ostensibly forecasting
(Table~\ref{tab:eq2_cpi_passthrough}).
The gap between these two responses is the forecast error documented
in Table~\ref{tab:eq3_forecast_error}.

% -------------------------------------------------------------
\subsection{Visualising the Treatment Variation}%
\label{sec:event_study}
% -------------------------------------------------------------

Figure~\ref{fig:meat_cpi_spaghetti} displays province-level
meat CPI trajectories over the sample period.

\begin{figure}[htbp]
\centering
\includegraphics[width=0.84\textwidth]{fig_meat_cpi_spaghetti.pdf}
\caption{Province-level meat CPI trajectories, 2016--2025.
  Each line is one province; colour indicates macro-region.
  The vertical dashed line marks August 2018 (first ASF case).
  The sharp divergence in 2019--2020 followed by mean reversion
  illustrates the temporary supply-shock nature of the ASF
  episode.}
\label{fig:meat_cpi_spaghetti}
\end{figure}

Three features are visible.
All 31 provinces experience a sharp spike in 2019--2020, but with
meaningful cross-province heterogeneity in magnitude.
The shock reverses by 2021--2022 as pork supply recovered.
The pre-2018 period shows moderate variation that provides
additional identifying power beyond the ASF episode.

\begin{figure}[htbp]
\centering
\includegraphics[width=0.80\textwidth]{fig_meat_shock_distribution.pdf}
\caption{Distribution of province-level meat price shocks
  by CFPS wave.  Each dot is one province; boxes show
  interquartile range.  The 2020 wave captures the ASF
  peak; the 2022 wave shows the reversal.}
\label{fig:meat_shock_dist}
\end{figure}

Figure~\ref{fig:meat_vs_headline} plots the meat shock against
forward realised headline CPI at the province level, providing
a direct visual counterpart to the regression in
Table~\ref{tab:eq2_cpi_passthrough}.
The negative slope confirms that high-shock provinces did not
experience higher subsequent headline inflation.

\begin{figure}[htbp]
\centering
\includegraphics[width=0.80\textwidth]{fig_meat_vs_headline.pdf}
\caption{Meat shock vs.\ forward realised headline CPI
  at the province level.
  The fitted line shows the negative relationship: provinces
  with larger meat shocks do not experience higher subsequent
  headline inflation.}
\label{fig:meat_vs_headline}
\end{figure}
